	A incerteza é um desafio grande na vida dos gestores, sejam eles públicos ou privados, e o combate a incerteza se da pela busca de informações, o risco é inversamente proporcional ao conhecimento que temos sobre determinado assunto,e o conhecimento só se da pela interpretação das informações obtidas através de um processo de mineração de dados, quando falamos de poucos dados, os modelos tradicionais de previsão conseguem acertar com um determinado nível de precisão, mas quando queremos resultados mais precisos, ou que nossos dados são mais complexos, evento cada vez mais comum na era do \textit{big-data}, os modelos tradicionais de previsão encontram limitações, nesse sentido a presente monografia busca responder em que medida modelos preditivos baseados em \ac{ia} conseguem estimar o PIB brasileiro com maior precisão do que os modelos tradicionais VAR, ARIMA, DFM e MIDAS? Para isso se utilizará de métodos quantitativos visando estimar a acurácia, custo computacional e interpretabilidade dos modelos testados, para enfim determinar se modelos de IA conseguem reduzir a incerteza.
	
	\palavraschave{Inteligência Artificial; PIB brasileiro; Métodos Quantitativos; Testes de modelos.}